\fontsize{13pt}{22.5pt}\selectfont

% ═══════════════════════════════════════════════
% Tiêu đề căn giữa, in hoa đậm giống DANH MỤC TỪ VIẾT TẮT (xem sec_abbrev.tex)
\phantomsection
\addcontentsline{toc}{section}{PHỤ LỤC}
\begin{center}
    {\large\bfseries PHỤ LỤC}
\end{center}


% ───────────────────────────────────────────────
\subsection*{Phụ lục A. Cấu hình huấn luyện PPO}
\addcontentsline{toc}{subsection}{Phụ lục A. Cấu hình huấn luyện PPO}

Dưới đây là tệp cấu hình YAML của PPO cho môi trường Football Table có kết hợp self-play,
đúng bản đã dùng cho lần chạy \texttt{FB01}. Những môi trường khác giữ cấu trúc tương tự,
chỉ chỉnh siêu tham số theo Bảng~\ref{tab:hyperparams}.
Cấu hình MAPPO hoàn toàn tương tự PPO và MA-POCA trên cùng một môi trường, chỉ khác ở
việc khai báo \texttt{trainer\_type: mappo}; vì vậy đề tài không liệt kê riêng, và mọi
chênh lệch kết quả giữa ba thuật toán này đều quy được về bản thân thuật toán.

\begin{lstlisting}[basicstyle=\ttfamily\scriptsize,caption={football\_ppo.yaml},captionpos=b]
behaviors:
  Football:
    trainer_type: ppo
    hyperparameters:
      batch_size: 2048
      buffer_size: 40960
      learning_rate: 3.0e-4
      beta: 0.005
      epsilon: 0.2
      lambd: 0.95
      num_epoch: 5
      learning_rate_schedule: linear
    network_settings:
      normalize: true
      hidden_units: 256
      num_layers: 2
    reward_signals:
      extrinsic:
        gamma: 0.99
        strength: 1.0
    self_play:
      save_steps: 50000
      team_change: 100000
      swap_steps: 200000
      window: 10
      play_against_latest_model_ratio: 0.5
      initial_elo: 1200.0
    keep_checkpoints: 10
    max_steps: 1600000
    time_horizon: 1000
    summary_freq: 10000
\end{lstlisting}

% ═══════════════════════════════════════════════
\subsection*{Phụ lục B. Cấu hình huấn luyện SAC}
\addcontentsline{toc}{subsection}{Phụ lục B. Cấu hình huấn luyện SAC}

Tệp cấu hình YAML của SAC cho môi trường Cross The Road, đúng như đã dùng cho lần chạy
\texttt{CrossTheRoad\_SAC\_01}. Tham số \texttt{steps\_per\_update} bằng $20$ là điểm cần chú
ý nhất: nó quyết định tốc độ hội tụ đo được của SAC, như phân tích ở
Chương~\ref{ch:results}.

\begin{lstlisting}[basicstyle=\ttfamily\scriptsize,caption={ctr\_sac.yaml},captionpos=b]
behaviors:
  CrossTheRoad:
    trainer_type: sac
    hyperparameters:
      batch_size: 256
      buffer_size: 100000
      learning_rate: 3.0e-4
      learning_rate_schedule: constant
      buffer_init_steps: 1000
      tau: 0.005
      steps_per_update: 20
      save_replay_buffer: false
      init_entcoef: 0.5
      reward_signal_steps_per_update: 20
    network_settings:
      normalize: false
      hidden_units: 128
      num_layers: 2
    reward_signals:
      extrinsic:
        gamma: 0.99
        strength: 1.0
      curiosity:
        gamma: 0.99
        strength: 0.02
        learning_rate: 3.0e-4
    keep_checkpoints: 5
    max_steps: 2000000
    time_horizon: 64
    summary_freq: 10000
\end{lstlisting}

Trên môi trường đối kháng, SAC vẫn dùng chính trainer \texttt{sac} của ML-Agents kết hợp
khối \texttt{self\_play} giống hệt Phụ lục A, nên lần chạy \texttt{Football\_SAC\_01} có
đầy đủ chỉ số ELO. Bên cạnh đó, đề tài cũng dựng sẵn con đường kết nối môi trường Unity
với một thư viện học tăng cường bên ngoài qua Python Low-Level API (xem
Mục~\ref{sub:pyapi}). Đoạn mã dưới đây minh họa lớp bọc môi trường; cần nhấn mạnh rằng
đây là phương án dự phòng đã kiểm chứng, không phải nguồn của bất kỳ số liệu nào trong
Chương~\ref{ch:results}.

\begin{lstlisting}[language=Python,caption={Kết nối môi trường Unity với mlagents\_envs + Stable-Baselines3},captionpos=b]
from mlagents_envs.environment import UnityEnvironment
from mlagents_envs.envs.unity_gym_env import UnityToGymWrapper
from stable_baselines3 import SAC

# Ket noi truc tiep voi moi truong Unity dang chay
unity_env = UnityEnvironment(file_name=None, base_port=5005)
env = UnityToGymWrapper(unity_env, allow_multiple_obs=False)

# Khoi tao tac nhan SAC voi mang 3 lop x 512 don vi
model = SAC(
    "MlpPolicy", env,
    learning_rate=3e-4,
    buffer_size=500000,
    batch_size=1024,
    tau=0.005,
    gamma=0.99,
    ent_coef="auto",
    policy_kwargs=dict(net_arch=[512, 512, 512]),
    tensorboard_log="./results/football_sac_v1",
    verbose=1,
)
model.learn(total_timesteps=10_000_000)
model.save("results/football_sac_v1/Football_sac")
env.close()
\end{lstlisting}

% ═══════════════════════════════════════════════
\subsection*{Phụ lục C. Cấu hình huấn luyện MA-POCA}
\addcontentsline{toc}{subsection}{Phụ lục C. Cấu hình huấn luyện MA-POCA}

Tệp cấu hình YAML của MA-POCA cho môi trường Capture The Flag, đúng như đã dùng cho lần
chạy \texttt{ctf\_poca\_v2}. Cấu hình MAPPO của lần chạy \texttt{ctf\_mappo\_v2} trùng
khớp từng tham số với tệp này, chỉ khác \texttt{trainer\_type: mappo}.

\begin{lstlisting}[basicstyle=\ttfamily\scriptsize,caption={ctf\_poca.yaml},captionpos=b]
behaviors:
  PuzzleBehavior:
    trainer_type: poca
    hyperparameters:
      batch_size: 1024
      buffer_size: 20480
      learning_rate: 3.0e-4
      beta: 0.01
      epsilon: 0.2
      lambd: 0.95
      num_epoch: 3
      learning_rate_schedule: linear
    network_settings:
      normalize: false
      hidden_units: 256
      num_layers: 3
      vis_encode_type: simple
    reward_signals:
      extrinsic:
        gamma: 0.99
        strength: 1.0
      curiosity:
        gamma: 0.99
        strength: 0.02
        learning_rate: 3.0e-4
    keep_checkpoints: 10
    max_steps: 1700000
    time_horizon: 128
    summary_freq: 20000
\end{lstlisting}

% ═══════════════════════════════════════════════
\subsection*{Phụ lục D. Cấu hình huấn luyện DQN}
\addcontentsline{toc}{subsection}{Phụ lục D. Cấu hình huấn luyện DQN}

Tệp cấu hình YAML của DQN cho môi trường Cross The Road, đúng như đã dùng cho lần chạy
\texttt{ctr\_dqn\_v1}. DQN được cài đặt dưới dạng trainer mở rộng (Mục~\ref{sub:pyapi})
nên khai báo \texttt{trainer\_type: dqn} chỉ hoạt động khi gói mở rộng đã được cài đặt.
Nhóm tham số \texttt{buffer\_size}, \texttt{tau} và \texttt{steps\_per\_update} được đặt
trùng với cấu hình SAC ở Phụ lục B, để hai thuật toán off-policy tiêu thụ dữ liệu theo
cùng một nhịp. Nhóm tham số \texttt{exploration\_*} là phần riêng của DQN, mô tả lịch
giảm tuyến tính của hệ số $\varepsilon$ trong quy tắc chọn hành động tham lam có nhiễu.

\begin{lstlisting}[basicstyle=\ttfamily\scriptsize,caption={ctr\_dqn.yaml},captionpos=b]
behaviors:
  CrossTheRoad:
    trainer_type: dqn
    hyperparameters:
      batch_size: 256
      buffer_size: 100000
      buffer_init_steps: 1000
      learning_rate: 3.0e-4
      learning_rate_schedule: constant
      tau: 0.005
      steps_per_update: 20
      save_replay_buffer: false
      exploration_schedule: linear
      exploration_initial_eps: 1.0
      exploration_final_eps: 0.05
      exploration_decay_steps: 500000
      reward_signal_steps_per_update: 20
    network_settings:
      normalize: false
      hidden_units: 128
      num_layers: 2
    reward_signals:
      extrinsic:
        gamma: 0.99
        strength: 1.0
      curiosity:
        gamma: 0.99
        strength: 0.02
        learning_rate: 3.0e-4
    keep_checkpoints: 5
    max_steps: 2000000
    time_horizon: 64
    summary_freq: 10000
\end{lstlisting}

Trên môi trường Football Table, tệp cấu hình của DQN giữ nguyên cấu trúc trên nhưng
\textbf{không có} khối \texttt{self\_play}: khi khai báo khối này, ML-Agents phát sinh
lỗi tra cứu chính sách trong bộ điều phối đối kháng và dừng ngay khi khởi tạo. Lần chạy
\texttt{football\_dqn\_v1} vì vậy dùng bản dựng \texttt{FootballDiscrete} với tham số
hành vi chuyển từ tám kênh liên tục sang tám nhánh ba mức, và không sinh ra chỉ số ELO.

% ═══════════════════════════════════════════════
\subsection*{Phụ lục E. Mã nguồn minh họa}
\addcontentsline{toc}{subsection}{Phụ lục E. Mã nguồn minh họa}

\noindent\textbf{E.1. Không gian quan sát (Cross The Road):}
\begin{lstlisting}[language=Java,caption={CollectObservations - CrossTheRoadAgent.cs},captionpos=b]
public override void CollectObservations(VectorSensor sensor)
{
    // 3 gia tri: toa do x, y, z cua agent
    sensor.AddObservation(transform.localPosition);
    // 3 gia tri: toa do x, y, z cua dich
    sensor.AddObservation(goal.transform.localPosition);
}
\end{lstlisting}

\noindent\textbf{E.2. Hàm phần thưởng (Cross The Road):}
\begin{lstlisting}[language=Java,caption={GivePoints / TakeAwayPoints - CrossTheRoadAgent.cs},captionpos=b]
public void GivePoints()          // Khi den dich
{
    AddReward(1.0f);
    EndEpisode();
}

public void TakeAwayPoints()      // Khi va cham phuong tien
{
    AddReward(-0.025f);
    EndEpisode();
}
\end{lstlisting}

\noindent\textbf{E.3. Phần thưởng nhóm (môi trường Capture The Flag):}
\begin{lstlisting}[language=Java,caption={Group reward - EnvController.cs (trich doan)},captionpos=b]
// Ca hai agent cung roi phong -> thuong nhom
if (agents[0].agent.ThisAgentLeft && agents[1].agent.ThisAgentLeft)
    agentGroup.AddGroupReward(0.5f / MaxEnvironmentSteps);

// Phat thoi gian (hurry-up penalty)
agentGroup.AddGroupReward(-0.25f / MaxEnvironmentSteps);

// Ca hai den checkpoint -> thuong lon + ket thuc episode
if (allFound) {
    agentGroup.AddGroupReward(reward);
    agentGroup.EndGroupEpisode();
}
\end{lstlisting}

\noindent\textbf{E.4. Phần thưởng định hướng (Football Table):}
\begin{lstlisting}[language=Java,caption={AddShotReward - FootballAgent.cs},captionpos=b]
void AddShotReward()
{
    if (Stats.HasBall) {
        Vector3 delta = opponentTeam.Goal.transform.position
                        - ball.transform.position;
        float reward = shotRewardMultiplier *
                       Vector3.Dot(delta.normalized, ball.Velocity);
        AddReward(reward);
    }
}
\end{lstlisting}

% ═══════════════════════════════════════════════
\subsection*{Phụ lục F. Kết quả TensorBoard}
\addcontentsline{toc}{subsection}{Phụ lục F. Kết quả TensorBoard}

Toàn bộ số liệu huấn luyện (Cumulative Reward, Episode Length, Policy Loss, Entropy,
ELO\ldots) đều được ghi lại tự động qua TensorBoard cho từng lần chạy và lưu trong
thư mục \texttt{results/} của dự án. Bảng~\ref{tab:tb_summary} tóm tắt các chỉ số
cuối cùng thực đo được của những lần chạy đã huấn luyện đầy đủ.

\begin{table}[H]
\centering
\caption{Kết quả thực đo của $15$ lần huấn luyện}
\label{tab:tb_summary}
\renewcommand{\arraystretch}{1.4}
\tablefont
\setlength{\tabcolsep}{4pt}
\begin{tabular}{|p{4.2cm}|>{\centering\arraybackslash}p{2.1cm}
                |>{\centering\arraybackslash}p{1.8cm}
                |>{\centering\arraybackslash}p{2.6cm}
                |>{\centering\arraybackslash}p{2.1cm}|}
\hline
\textbf{Lần chạy} & \textbf{Thuật toán} & \textbf{Số bước} & \textbf{Reward TB (10\% cuối)} & \textbf{Thời gian} \\
\hline
\multicolumn{5}{|l|}{\textbf{Cross The Road}} \\
\hline
ctr01 & PPO & 2M & $0{,}690$ & $\sim$48 phút \\
\hline
CrossTheRoad\_SAC\_01 & SAC & 2M & $0{,}846$ & $\sim$139 phút \\
\hline
ctr\_poca\_v1 & MA-POCA & 2M & $0{,}556$ & $\sim$59 phút \\
\hline
ctr\_mappo\_v1 & MAPPO & 2M & $0{,}835$ & $\sim$58 phút \\
\hline
ctr\_dqn\_v1 & DQN & 2M & $0{,}336$ & $\sim$148 phút \\
\hline
\multicolumn{5}{|l|}{\textbf{Capture The Flag}} \\
\hline
ctf\_ppo\_v1 & PPO & 1{,}70M & $0{,}834$ & $\sim$50 phút \\
\hline
ctf\_sac\_v1$^{\dagger}$ & SAC & 1{,}70M & $0{,}183$ & $\sim$152 phút \\
\hline
ctf\_poca\_v2 & MA-POCA & 1{,}70M & $0{,}838$ & $\sim$53 phút \\
\hline
ctf\_mappo\_v2 & MAPPO & 1{,}70M & $0{,}840$ & $\sim$52 phút \\
\hline
ctf\_dqn\_v1 & DQN & 1{,}70M & $0{,}800$ & $\sim$258 phút \\
\hline
\multicolumn{5}{|l|}{\textbf{Football Table}} \\
\hline
FB01 & PPO & 1{,}60M & $1{,}691$ \newline (ELO 1259) & $\sim$47 phút \\
\hline
Football\_SAC\_01 & SAC & 1{,}60M & $1{,}842$ \newline (ELO 1254) & $\sim$165 phút \\
\hline
Football\_POCA\_01 & MA-POCA & 1{,}60M & $1{,}752$ \newline (ELO 1237) & $\sim$50 phút \\
\hline
Football\_MAPPO\_01 & MAPPO & 1{,}60M & $2{,}117$ \newline (ELO 1260) & $\sim$48 phút \\
\hline
football\_dqn\_v1$^{\ddagger}$ & DQN & 1{,}60M & $1{,}825$ \newline (không có ELO) & $\sim$147 phút \\
\hline
\end{tabular}

\vspace{0.2cm}
{\footnotesize
$^{\dagger}$ Thời gian của \texttt{ctf\_sac\_v1} chỉ tính đoạn $1{,}0$M--$1{,}7$M sau khi khôi phục từ checkpoint; đoạn $0$--$1{,}0$M chạy trong một phiên trước đó bị ngắt nên không đo được liền mạch.

\smallskip
$^{\ddagger}$ Riêng DQN trên Football phải tắt self-play, do ML-Agents phát sinh lỗi khi kết hợp trainer off-policy với ghost trainer. Vì vậy lần chạy này không có chỉ số ELO và \textbf{không so sánh trực tiếp được} với bốn thuật toán còn lại trên cùng môi trường.

\smallskip
Ghi chú chung: mọi lần chạy đều dùng $8$ training area song song trong mỗi scene. Reward là
phần thưởng trung bình mười phần trăm cuối, cùng định nghĩa dùng trong Chương 4.}
\end{table}

\vspace{0.5cm}
\noindent\textbf{Mã nguồn và mô hình:} Toàn bộ mã nguồn dự án Unity, các tệp cấu hình,
script huấn luyện và các mô hình đã huấn luyện (định dạng \texttt{.onnx}) được lưu trữ
tại kho mã nguồn của đề tài.
